% ----------------------------------------------------------------
\subsection{International Film Industry}
\label{sec:film_industry}
% ----------------------------------------------------------------
A variety of factors influence the distribution of international film production including labour costs, government incentives, co-production treaties, and the existence of pre-existing production networks \citep{morawetz2007, owens2020, mitric2018}. In recent decades, a growing number of national and subnational governments have introduced schemes to attract international film productions \citep{christopherson2010}. These incentive schemes take different forms including cash rebates and tax credits, and vary substantially across jurisdictions with respect to criteria such as credit rates, eligibility criteria, and cultural content requirements \citep{olsberg2025}. 

The net economic impact of production incentives is debatable. The fiscal cost of film production subsidies often outstrips the tax revenue they generate \citep{luther2010, tannenwald2010}, and several studies have found evidence of ``crowding out'' where film production incentives merely draw resources from other sectors of the economy \citep{calcagno2004}. Conversely, many industry-commissioned studies have found substantial positive returns from incentive schemes \citep{ey2009, olsberg2023}. However, the methodological basis of these studies has been questioned \citep{button2021}.

A related but limited set of literature sets aside the issue of net welfare effects and examines whether incentives succeed at attracting film productions. These studies focus exclusively on incentives at the U.S.\ state level. This is due to the large number of state-level schemes available and their similar characteristics reducing potential confounding factors. Some of these studies find meaningful evidence that incentive schemes can attract film productions. \citet{button2021} employs synthetic control and difference-in-differences methods to evaluate the programmes in Louisiana and New Mexico (two early adopters of generous state film incentives) and finds that incentives lead to a statistically significant increase in the number of films produced. \citet{workman2021} exploits the fact that California's Film and Television Tax Credit Program allocates tax incentives by lottery, creating quasi-random variation in which productions receive incentives and allowing him to estimate causal effects more cleanly than other studies. He finds that the offer of an incentive meaningfully increased the number of films produced in the state.

Other studies are more mixed. \citet{owens2020} estimate a discrete choice model of filming location decisions for productions released between 1999 and 2013. They find significant heterogeneity in responsiveness: mid-sized studios respond to all forms of incentives, major studios respond only to refundable and transferable tax credits, and independent productions exhibit no statistically significant sensitivity to incentives. Finally, \citet{button2019} uses a panel difference-in-differences design comparing U.S.\ states before and after they adopt film incentives against similar states that did not adopt over the same period. He finds an increase in the production of TV series but no meaningful effect on feature film production.

% ----------------------------------------------------------------
\subsection{The Gravity Model of Trade}
\label{sec:gravity}
% ----------------------------------------------------------------

The gravity model of trade predicts that bilateral trade between two countries is proportional to the product of their economic sizes and inversely related to the geographic distance between them \citep{tinbergen1962}. It has been shown to be derivable from a wide range of popular trade models \citep{eaton2002, krugman1980, chaney2008}, and has been extensively used to model trade flows \citep{head2014}. When estimating gravity models of international trade, it has become standard to include dummy variables for each exporting and importing country in the sample to control for all unobservable country-specific factors that affect a country's overall propensity to trade \citep{anderson2003}.

The gravity model has been extensively applied to goods trade but employed less frequently to trade in services. \citet{santanna2023} attribute this to three factors: well-known limitations in bilateral services trade data, the comparatively recent emergence of services as a significant component of international trade, and a widespread assumption that insights from the goods trade literature transfer directly to services. However, existing literature suggests that the gravity framework is well suited to modelling services flows. \citet{kimura2006} estimate gravity equations for bilateral services trade between ten OECD member countries and a broader sample of trading partners. Their principal finding is that services trade is better predicted by the gravity specification than goods trade, and that there exists a complementary relationship between goods exports and services imports.

\citet{walsh2006} applies the gravity model to aggregate services trade and finds that the gravity model fits services trade in a manner broadly comparable to goods. However, he finds that GDP per capita and common language emerge as the most significant determinants of services trade, while geographic distance is generally insignificant. This is consistent with the observation that many services do not require physical shipment and are therefore less sensitive to transport costs but notably diverges from the goods gravity literature which finds distance to be the dominant friction in international trade. \citet{thai2006}, analysing Vietnam--EU bilateral services flows, confirms these patterns and documents the significance of GDP and common language as determinants.
